
\documentclass[a4paper,11pt]{article}
\usepackage[a4paper, margin=8em]{geometry}

% usa i pacchetti per la scrittura in italiano
\usepackage[french,italian]{babel}
\usepackage[T1]{fontenc}
\usepackage[utf8]{inputenc}
\frenchspacing 

% usa i pacchetti per la formattazione matematica
\usepackage{amsmath, amssymb, amsthm, amsfonts}

% usa altri pacchetti
\usepackage{gensymb}
\usepackage{hyperref}
\usepackage{standalone}

% imposta il titolo
\title{Appunti Comunicazioni Numeriche}
\author{Luca Seggiani}
\date{2026}

% imposta lo stile
% usa helvetica
\usepackage[scaled]{helvet}
% usa palatino
\usepackage{palatino}
% usa un font monospazio guardabile
\usepackage{lmodern}

\renewcommand{\rmdefault}{ppl}
\renewcommand{\sfdefault}{phv}
\renewcommand{\ttdefault}{lmtt}

% disponi il titolo
\makeatletter
\renewcommand{\maketitle} {
	\begin{center} 
		\begin{minipage}[t]{.8\textwidth}
			\textsf{\huge\bfseries \@title} 
		\end{minipage}%
		\begin{minipage}[t]{.2\textwidth}
			\raggedleft \vspace{-1.65em}
			\textsf{\small \@author} \vfill
			\textsf{\small \@date}
		\end{minipage}
		\par
	\end{center}

	\thispagestyle{empty}
	\pagestyle{fancy}
}
\makeatother

% disponi teoremi
\usepackage{tcolorbox}
\newtcolorbox[auto counter, number within=section]{theorem}[2][]{%
	colback=blue!10, 
	colframe=blue!40!black, 
	sharp corners=northwest,
	fonttitle=\sffamily\bfseries, 
	title=~\thetcbcounter: #2, 
	#1
}

% disponi definizioni
\newtcolorbox[auto counter, number within=section]{definition}[2][]{%
	colback=red!10,
	colframe=red!40!black,
	sharp corners=northwest,
	fonttitle=\sffamily\bfseries,
	title=~\thetcbcounter: #2,
	#1
}

% U.D.M
\newcommand{\amp}{\ensuremath{\mathrm{A}}}
\newcommand{\volt}{\ensuremath{\mathrm{V}}}
\newcommand{\meter}{\ensuremath{\mathrm{m}}}
\newcommand{\second}{\ensuremath{\mathrm{s}}}
\newcommand{\farad}{\ensuremath{\mathrm{F}}}
\newcommand{\henry}{\ensuremath{\mathrm{H}}}
\newcommand{\siemens}{\ensuremath{\mathrm{S}}}

% circuiti
\usepackage{circuitikz}
\usetikzlibrary{babel}

% disegni
\usepackage{pgfplots}
\pgfplotsset{width=10cm,compat=1.9}

% disponi codice
\usepackage{listings}
\usepackage[table]{xcolor}

\lstdefinestyle{codestyle}{
		backgroundcolor=\color{black!5}, 
		commentstyle=\color{codegreen},
		keywordstyle=\bfseries\color{magenta},
		numberstyle=\sffamily\tiny\color{black!60},
		stringstyle=\color{green!50!black},
		basicstyle=\ttfamily\footnotesize,
		breakatwhitespace=false,         
		breaklines=true,                 
		captionpos=b,                    
		keepspaces=true,                 
		numbers=left,                    
		numbersep=5pt,                  
		showspaces=false,                
		showstringspaces=false,
		showtabs=false,                  
		tabsize=2
}

\lstdefinestyle{shellstyle}{
		backgroundcolor=\color{black!5}, 
		basicstyle=\ttfamily\footnotesize\color{black}, 
		commentstyle=\color{black}, 
		keywordstyle=\color{black},
		numberstyle=\color{black!5},
		stringstyle=\color{black}, 
		showspaces=false,
		showstringspaces=false, 
		showtabs=false, 
		tabsize=2, 
		numbers=none, 
		breaklines=true
}

\lstdefinelanguage{javascript}{
	keywords={typeof, new, true, false, catch, function, return, null, catch, switch, var, if, in, while, do, else, case, break},
	keywordstyle=\color{blue}\bfseries,
	ndkeywords={class, export, boolean, throw, implements, import, this},
	ndkeywordstyle=\color{darkgray}\bfseries,
	identifierstyle=\color{black},
	sensitive=false,
	comment=[l]{//},
	morecomment=[s]{/*}{*/},
	commentstyle=\color{purple}\ttfamily,
	stringstyle=\color{red}\ttfamily,
	morestring=[b]',
	morestring=[b]"
}

% disponi sezioni
\usepackage{titlesec}

\titleformat{\section}
	{\sffamily\Large\bfseries} 
	{\thesection}{1em}{} 
\titleformat{\subsection}
	{\sffamily\large\bfseries}   
	{\thesubsection}{1em}{} 
\titleformat{\subsubsection}
	{\sffamily\normalsize\bfseries} 
	{\thesubsubsection}{1em}{}

% disponi alberi
\usepackage{forest}

\forestset{
	rectstyle/.style={
		for tree={rectangle,draw,font=\large\sffamily}
	},
	roundstyle/.style={
		for tree={circle,draw,font=\large}
	}
}

% disponi algoritmi
\usepackage{algorithm}
\usepackage{algorithmic}
\makeatletter
\renewcommand{\ALG@name}{Algoritmo}
\makeatother

% disponi numeri di pagina
\usepackage{fancyhdr}
\fancyhf{} 
\fancyfoot[L]{\sffamily{\thepage}}

\makeatletter
\fancyhead[L]{\raisebox{1ex}[0pt][0pt]{\sffamily{\@title \ \@date}}} 
\fancyhead[R]{\raisebox{1ex}[0pt][0pt]{\sffamily{\@author}}}
\makeatother

\begin{document}
% sezione (data)
\section{Lezione del 24-04-26}

% stili pagina
\thispagestyle{empty}
\pagestyle{fancy}

% testo
Concluso l'argomento del campionamento dei segnali, passiamo a discutere la \textit{teoria dei codici} per quanto riguarda la parte della codifica dei segnali per poi trasmetterli su un mezzo. 

\subsection{Teoria dei codici}
L'obiettivo della codifica di sorgente e di canale al trasmettitore è, rispettivamente:
\begin{itemize}
	\item Per quanto riguarda la \textit{codifica di sorgente} (\textbf{source encoder}), questa rimuove la ridondanza del segnale, per ottimizzare la quantità di informazione inviata (compressione);
	\item Per quanto riguarda la \textit{codifica di canale} (\textbf{channel encoder}), questa introduce nuova ridondanza, sotto forma di meccanismi di protezione dagli errori.
\end{itemize}

In ricezione, i decoder ricostruiscono:
\begin{itemize}
	\item Prima la parola di codice dal canale;
	\item Quindi il segnale originale dalla parola di codice (che codificava il segnale originale stesso).
\end{itemize}

Esistono 2 categorie principali di codici:
\begin{enumerate}
	\item I codici a \textbf{rivelazione} di errore individuano la presenza di errori, senza correggerli;
	\item I codici a \textbf{correzione} di errore riescono sia ad individuare che a correggere gli errori.
\end{enumerate}

\subsubsection{Codici di canale lineari}
I \textit{codici di canale lineari} sono particolari tipi di codici basati su combinazioni lineari di bit. Di questi ne esistono altre 2 categorie:
\begin{itemize}
	\item Codici a \textbf{blocco}: operano su blocchi finiti di bit;
	\item Codici a \textbf{convoluzione}: operano su sequenze con memoria (quindi stream) di bit.
\end{itemize}

Il \textit{rapporto} $R$ di un codice a blocco è il rapporto fra numero di bit in ingresso $k$ e numero di bit in uscita $n$ per unità di trasmissione:
$$
R = \frac{k}{n}
$$
Idealmente vorremo un rapporto il più basso possibile, in quanto questo significherà che introdurremo meno ridondanza possibile.
Il rapporto di un codice è anche legato al tempo di trasmissione di una parola. Infatti varrà:
$$
k T_b = n T_c
$$
dove $T_b$ è il tempo di ciascun bit di ingresso e $T_c$ il tempo di ciascun bit di uscita.
L'uguaglianza dovrà reggere affinché il sistema funzioni in tempo continuo, cioè si riescano a reintrodurre bit nel mezzo con la stessa frequenza con la quale si ottengono dal campionatore. 
Da questo seguirà che:
$$
T_c = \frac{k}{n} T_b \leq T_b
$$
cioè il tempo di trasmissione sul canale dovrà essere \textit{minore} del tempo di ricezione di bit di ingresso. Riguardo alle frequenze $R_c$ e $R_b$ di uscita e di ingresso rispettivamente, varrà invece:
$$
R_c = \frac{1}{T_c} = \frac{n}{k} \frac{1}{T_b} = \frac{n}{k}R_b \geq R_b
$$
cioè in maniera analoga a prima, la frequenza di trasmissione sul canale dovrà essere \textit{maggiore} della frequenza di ricezione di bit di ingresso. 

\subsubsection{Codici a blocco}
Iniziamo a vedere 2 semplici esempi di codice a blocco.

\begin{enumerate}
\item \textbf{\textsf{Codice a ripetizione}} \\
Il codice a ripetizionei è uno dei codici a blocco più semplici che esistano. Consiste nel riprodurre il bit che vogliamo trasmettere $n$ volte sul canale di sucita. Ad esempio, un codice di ripetizione con $R = \frac{1}{3}$ mappa i bit direttamente a:
$$
m = 0 \rightarrow c = [000], \quad
m = 1 \rightarrow c = [111]
$$

Potremmo quindi chiederci come decodificare il codice, tenendo presente che potrebbero esserci errori (bit-flip) lungo il segnale ricevuto. La risposta è scegliendo il bit più frequente nella parola ricevuta, cioè implementando la cosiddetta \textit{decodifica a maggioranza}.

Un codice a ripetizione può \textit{rivelare} fino ad $n - 1$ errori, in quanto si possono rivoltare fino a $n - 1$ prima di ottenere quella che è nuovamente una parola del codice (parole che contengono sia $0$ che $1$ sono automaticamente sbagliate per la definizione di codici a ripetizione).

Allo stesso modo, riesce a \textit{correggere} fino a $\frac{n - 1}{2}$ errori, in quanto quando oltre la metà dei bit viene rivoltata si può si rilevare che si è verificato un errore, ma non verificare se si è risolto in maniera corretta.

Il tradeoff dei codici a ripetizione è che si spende molto più tempo a trasmettere lo stesso bit, in particolare con un fattore $n$. Nel caso in cui vogliamo trasmettere con le ipotesi di prima, cioè con la stessa frequenza di informazione all'ingresso e all'uscita sul canale, l'unica cosa che possiamo fare è aumentare la banda all'uscita, cioè su cui trasmettiamo il segnale. 

Dal punto di vista probabilistico, data la probabilità di sbagliare un bit come $p$, la parola di sbagliare $t$ bit in una parola di $n$ bit è data dalla formula di Bernoulli:
$$
p(t, n) = \binom{n}{t} p^t (1 - p)^{n - t}
$$
dove:
$$
\binom{n}{t} = \frac{n!}{t!(n - t)!}
$$
Notiamo che nei sistemi di comunicazione reali la probabilità di sbagliare un bit è nell'ordine di $10^{-2} - 10^{-5}$.
Siccome siamo capaci di correggere un errore di al massimo $\frac{n - 1}{2}$ bit, la probabilità di errore della parola di codice è pari a:
$$
P = \sum_{t = \frac{n + 1}{2}}^n \binom{n}{t} p^t (1 - p)^{n - t} \approx 
\binom{n}{\frac{n + 1}{2}} p^{\frac{n + 1}{2}} (1 - p)^{\frac{n - 1}{2}} \approx \binom{n}{\frac{n + 1}{2}} p^{\frac{n + 1}{2}}
$$
dove facciamo l'approssimazione di porre la maggior pobabilità sull'evento di errore più probabile.

\item \textbf{\textsf{Codice a controllo di parità}} \\
Consideriamo un altro semplice tipo di codice di rivelazione degli errori, con rapporto:
$$
R = \frac{k}{k + 1}
$$
Abbiamo che di questi, usiamo:
\begin{itemize}
	\item $k$ bit di informazione effettivamente utile;
	\item $k + 1$ bit di informazione inviata, di cui il bit in più è un bit di \textit{parità} del blocco utile. Questo viene calcolato facendo la somma modulo 2 dei bit nella stringa e trascurando il riporto.
\end{itemize}

Senza bit di parità, un solo bit sbagliato fa sbagliare tutta la parola. Riprendendo la formula di probabiltà si ha:
$$
P = \sum_{t = 1}^{k} \binom{k}{t} p^t (1 - p)^{k - t}
$$
che è piuttosto grande. Se aggiungo un bit di parità, la parola diventa di 12 bit, e si sbaglia quando si fanno almeno 2 errori. Gli errori di 1 bit vengono rivelati e corretti (mediante ritrasmissione). Notiamo che dal fatto che la maggior parte di errori si presentano come \textit{spray}, la probabilità di un errore di 1 bit è circa del $50 \%$. Riprendiamo quindi la formula di probabilità, stavolta:
$$
P = \sum_{t = 2}^{k} \binom{k}{t} p^t (1 - p)^{k - t}
$$
che è molto meglio.

\end{enumerate}

\subsubsection{Codici a blocco}
I \textbf{codici a blocco} sono rappresentati da trasformazioni lineari di stringhe di $k$ bit, dette \textit{blocchi} $\mathbf{b}$:
$$
\mathbf{b} = [ b_1, ..., b_k ]
$$
Definendo il codice a blocco in maniera tabulare, dove associamo ad ogni blocco di ingressa un blocco di uscita di lunghezza $n$, un codice a blocco lineare sarà l'\textit{insieme} delle $2^k$ parole $\mathbf{c}$ di uscita generate dalla trasformazione lineare della stringa $\mathbf{b}$:
$$
\mathbf{c} = [c_1, ..., c_n] = \mathbf{b} G = \sum_{i = 1}^k b_i \mathbf{g}_i
$$
dove $G$ è una matrice $k \times n$ detta \textit{matrice generatrice} del codice. Le $\mathbf{g}_i$ a questo punto sono le $k$ righe di dimensione $n$ della matrice. 

Le proprietà dei codici a blocco (facilmente verificabili) sono:
\begin{enumerate}
	\item 
Ogni parola di codice è combinazione lineare delle righe di $G$;
	\item
Il codice è costituito da tutte le possibili combinazioni delle righe di $G$;
	\item
La somma di due parole di codice è ancora una parola di codice;
	\item
La n-upla di tutti zeri è sempre una parola di codice.
\end{enumerate}

\subsubsection{Distanza di Hamming}
Definiamo \textit{distanza di Hammming} fra due parole di codice $\mathbf{c}_1$ e $\mathbf{c}_2$ come il numero di posizioni di cui differiscono le parole di codice. La indichiamo come $d_H$:
$$
d_H(\mathbf{c}_1, \mathbf{c}_2)
$$

Il \textit{peso di Haming} di una parola di codice $\mathbf{c}$, a questo punto, è definito come la distanza dall'origine della parola:
$$
w(\mathbf{c}) = d_H(\mathbf{c}, 0)
$$
In altre parole, il peso è il numero di bit a 0 di una parola $\mathbf{c}$.

A questo punto, la \textit{distanza minima} di un codice è la minima distanza di Hamming fra tutte le parole di codice:
$$
d_{min} = \mathrm{min}_{i, j} d_H(\mathbf{c}_i,, \mathbf{c}_j) = \mathrm{min}_k d_H(\mathbf{c}_k, 0), \quad i \neq j
$$

\end{document}
